%%%PAGE 61 rot=0 legibility=good summary=线框穿磁场的 i-t 图像规律（单场匀速、反向场匀速、单一磁场匀加速、反向场匀加速、菱形线框）
%%%BLOCK id=061-01 ch=12 sec=线框穿磁场·图像 kind=summary alt=none cont=none y=0.02-0.14
\textbf{线框穿磁场}\quad 考察右手定则、电源串联、有效长度 % CHECK: 原稿“穿”字小写在“传”字上方（疑为更正），主行写作“线框传磁场”
%FIG p061_a 0.39 0.03 0.57 0.068
\notefig[0.3]{figures/p061_a.png}
% 图中文字：相加、相减（两组电池串联符号）
\textbf{线框图像}\quad 同形状
%%%END
%%%BLOCK id=061-02 ch=12 sec=线框穿磁场·图像 kind=method alt=none cont=none y=0.14-0.31
\textbf{单场匀速}\quad 进出反向轴对称
%FIG p061_b 0.10 0.173 0.25 0.30
\notefig[0.3]{figures/p061_b.png}
% 图中文字：θ、v、sin x 形、L
\[ E=B\cdot v\tan\theta\,t\cdot v \]
%FIG p061_c 0.458 0.17 0.60 0.245
\notefig[0.3]{figures/p061_c.png}
%FIG p061_d 0.265 0.243 0.41 0.306
\notefig[0.3]{figures/p061_d.png}
% 图中文字：i、t、2L/v_0
\[ Q=\left(\frac{\dfrac{BLv}{\sqrt{2}}}{R}\right)^{2}\frac{2L}{v_0} \] % CHECK: 按原稿转录，式中未见电阻 R 因子（Q=I^2Rt 应有 R）
%%%END
%%%BLOCK id=061-03 ch=12 sec=线框穿磁场·图像 kind=method alt=none cont=none y=0.31-0.49
\textbf{反向场，两同夹异\ 中加倍}
%FIG p061_e 0.10 0.32 0.56 0.495
\notefig[0.55]{figures/p061_e.png}
% 图中文字：匀速、i、t（两个 i-t 图像）
%%%END
%%%BLOCK id=061-04 ch=12 sec=线框穿磁场·图像 kind=method alt=none cont=none y=0.51-0.62
\textbf{单一磁场匀加速}
%FIG p061_f 0.20 0.538 0.32 0.604
\notefig[0.25]{figures/p061_f.png}
%FIG p061_j 0.335 0.527 0.49 0.611
\notefig[0.3]{figures/p061_j.png}
% 图中文字：a、i、t
前后对称后共线
%%%END
%%%BLOCK id=061-05 ch=12 sec=线框穿磁场·图像 kind=method alt=none cont=none y=0.61-0.74
\textbf{反向场匀加速}
%FIG p061_g 0.24 0.632 0.335 0.70
\notefig[0.3]{figures/p061_g.png}
%FIG p061_h 0.34 0.612 0.482 0.74
\notefig[0.3]{figures/p061_h.png}
两同夹异中加倍，（前后共线\ 中加倍）\\
前后图象共有线
%%%END
%%%BLOCK id=061-06 ch=12 sec=线框穿磁场·图像 kind=method alt=none cont=none y=0.74-0.87
%FIG p061_i 0.13 0.74 0.60 0.86
\notefig[0.6]{figures/p061_i.png}
% 图中文字：v→、L、2、i、t、进、过 L/2、未到界线之已进、往外出之往里进（图中有一处被修正带涂白）
图中批注：「进」「过$\dfrac{L}{2}$」；\unclear{未到界线}之已进；往外出之往里进 % CHECK: “未到界线”“往外”处字迹难辨
%%%END
%%%PAGEEND 61
